% from ega3-1-fr.tex
\oldpage[III]{156}

que $\widehat g_*\bigl(\widehat g^{,*}(\mathcal F)\bigr)$ est un
$\mathcal O_{\widehat X}$-Module algébrisable
(\hyperref[III.5.3.3-fr]{5.3.3}) puisque $g$ est propre. On peut maintenant
appliquer à $\mathcal F$ et à $g$ le résultat de
(\hyperref[III.5.3.4-fr]{5.3.4})~: le noyau $\mathcal P$ et le conoyau
$\mathcal R$ de l'homomorphisme
$\rho_{\mathcal F}:\mathcal F\to
\widehat g_*\bigl(\widehat g^{,*}(\mathcal F)\bigr)$ sont annulés par une
puissance $\widehat{\mathcal M}^{,n}$~; soient $T$ le sous-préschéma fermé de $X$
défini par $\mathcal M^n$, ayant $X-U$ pour espace sous-jacent,
$j:T\to X$ l'injection canonique, de sorte que le prolongement aux complétés
$\widehat j:\widehat T\to\widehat X$ est l'injection canonique (I, 10.14.7). On
peut donc écrire
$\mathcal P=\widehat j_*\bigl(\widehat j^{,*}(\mathcal P)\bigr)$ et
$\mathcal R=\widehat j_*\bigl(\widehat j^{,*}(\mathcal R)\bigr)$ et comme $U$
n'est pas vide, il résulte de l'hypothèse de récurrence que
$\widehat j^{,*}(\mathcal P)$ et $\widehat j^{,*}(\mathcal R)$ sont des
$\mathcal O_{\widehat T}$-Modules algébrisables~; en vertu de
(\hyperref[III.5.3.3-fr]{5.3.3}), $\mathcal P$ et $\mathcal R$ sont
algébrisables, et on peut alors appliquer (\hyperref[III.5.3.2-fr]{5.3.2}), qui
prouve finalement que $\mathcal F$ est algébrisable.

\hfill C.Q.F.D.
\end{env}

\subsection{Application : comparaison de morphismes de schémas usuels et de morphismes de schémas formels. Schémas formels algébrisables}
\label{subsection:III.5.4-fr}

\begin{theorem}[5.4.1]
\label{III.5.4.1-fr}
Soient $A$ un anneau adique noethérien, $\mathfrak J$ un idéal de définition de
$A$, $S=\operatorname{Spec}(A)$, $S'=\mathrm V(\mathfrak J)$. Soient
$u:X\to S$ un morphisme propre, $v:Y\to S$ un morphisme séparé de type fini, et
soient $\widehat S$, $\widehat X$, $\widehat Y$ les complétés de $S$, $X$, $Y$ le
long de $S'$, $u^{-1}(S')$, $v^{-1}(S')$ respectivement. Si, pour tout
$S$-morphisme $f:X\to Y$, $\widehat f:\widehat X\to\widehat Y$ est le
prolongement de $f$ aux complétés, l'application $f\mapsto\widehat f$ est une
bijection
\[
\operatorname{Hom}_S(X,Y)\xrightarrow{\ \sim\ }
\operatorname{Hom}_{\widehat S}(\widehat X,\widehat Y).
\]
\end{theorem}

\begin{proof}
Montrons d'abord que $f\mapsto\widehat f$ est injective. Supposons en effet que deux
$S$-morphismes $f$, $g$ de $X$ dans $Y$ soient tels que $\widehat f=\widehat g$.
On sait alors (I, 10.9.4) qu'il existe un voisinage ouvert $V$ de
$X'=u^{-1}(S')$ dans lequel $f$ et $g$ coïncident. Or, comme $u$ est une application
fermée, on a $V=X$ (\hyperref[III.5.1.3.1-fr]{5.1.3.1}), d'où $f=g$.

Prouvons maintenant que $f\mapsto\widehat f$ est surjective, et soit donc $h$ un
$\widehat S$-morphisme $\widehat X\to\widehat Y$. Soit $Z=X\times_S Y$, et
désignons par $p:Z\to X$ et $q:Z\to Y$ les projections canoniques~; $Z$ est de
type fini sur $S$ (I, 6.3.4), donc noethérien~; désignons par $\widehat Z$ son
complété le long de $Z'=p^{-1}(u^{-1}(S'))$~; on sait que $\widehat Z$ s'identifie
canoniquement à $\widehat X\times_{\widehat S}\widehat Y$, les projections
$\widehat Z\to\widehat X$ et $\widehat Z\to\widehat Y$ s'identifiant aux
prolongements $\widehat p$ et $\widehat q$ (I, 10.9.7). Comme $Y$ est séparé sur
$S$, $\widehat Y$ est séparé sur $\widehat S$ (I, 10.15.7), donc le morphisme
graphe
\[
\Gamma_h=(1_{\widehat X},h):\widehat X\to\widehat Z
\]
est une immersion fermée (I, 10.15.4). Soient $\mathfrak T$ le sous-préschéma
formel fermé de $\widehat Z$ associé à cette immersion, et
$j:\mathfrak T\to\widehat Z$ l'injection canonique, de sorte que
$\Gamma_h=j\circ w$, où $w:\widehat X\to\mathfrak T$ est un isomorphisme
(I, 10.14.3) dont l'isomorphisme réciproque est $\widehat p\circ j$~; en outre,
$\mathfrak T$ est évidemment propre sur $\widehat S$, puisque $\widehat X$ l'est~;
on en conclut (\hyperref[III.5.1.8-fr]{5.1.8}) qu'il existe un sous-préschéma
fermé $T$ de $Z$ tel que
$\mathfrak T=\widehat T=T_{/(T\cap Z')}$, et que $j=\widehat i$, où $i$ est
l'injection canonique $T\to Z$ (I, 10.14.7). Alors $p\circ i:T\to X$ est un
isomorphisme, car il en est ainsi de
$\widehat{p\circ i}=\widehat p\circ\widehat i$ par hypothèse, et il suffit
d'appliquer

